% Apunts de teoria en PDF · 4t ESO · Unitat 3: Equacions i inequacions
% Mateix estil que els apunts de Batxillerat (franja i caixes de teoria en blau, exemples en roig), però més visuals,
% curts i amb dibuixos, per a l'alumnat de 3r ESO. Es compila amb pdflatex dues vegades (per als enllaços de l'índex),
% amb logo-fentmates.pdf (_eines/logo/pdf/horitzontal-clar.pdf) al costat.
\documentclass[11pt,a4paper]{article}
\usepackage[utf8]{inputenc}
\usepackage[T1]{fontenc}
\usepackage[catalan]{babel}
\usepackage[sfdefault,book,lining]{FiraSans}
\usepackage[italic,noplusnominus,noequal]{mathastext}
\linespread{1.18}
\usepackage[a4paper,top=2.4cm,bottom=2.2cm,left=1.6cm,right=1.6cm,headheight=1.1cm,headsep=0.5cm,footskip=1cm]{geometry}
\usepackage{amsmath,amssymb,cancel}
\usepackage{icomma}
\usepackage{graphicx}
\usepackage[table]{xcolor}
\usepackage{tikz}\usetikzlibrary{calc,positioning,arrows.meta,decorations.pathreplacing,shapes.geometric}
\usepackage{fontawesome5}
\usepackage{ccicons}
\usepackage{fancyhdr}
\usepackage{enumitem}
\usepackage{array}
\usepackage{needspace}
\usepackage[most]{tcolorbox}
\usepackage{titlesec}
\usepackage[hidelinks]{hyperref}
\hypersetup{pdftitle={Apunts de teoria. Unitat 3: Equacions i inequacions. 4t ESO},pdfsubject={Matemàtiques B 4t ESO},pdfkeywords={equacions de segon grau, biquadrades, Ruffini, fraccions algebraiques, radicals, exponencials, logarítmiques, inequacions, taula de signes}}

% colors de la web (bloc d'Àlgebra: blau)
\definecolor{fmnavy}{HTML}{1B2748}
\definecolor{fmblue}{HTML}{3D6AA0}
\definecolor{fmbluel}{HTML}{4E81BD}
\definecolor{fmyellow}{HTML}{FFC94D}
\definecolor{fmgold}{HTML}{E9A919}
\definecolor{fmaccent}{HTML}{FF7A50}
\definecolor{fmdef}{HTML}{1D4F73}\definecolor{fmdefbg}{HTML}{EEF4F9}
\definecolor{fmex}{HTML}{A3203A}\definecolor{fmexbg}{HTML}{FBF0F2}
\definecolor{fmsf}{HTML}{C2410C}
\definecolor{fmok}{HTML}{15803D}
\definecolor{fmmuted}{HTML}{5D6685}
\definecolor{fmsurface}{HTML}{EEF1F8}
\definecolor{fmA}{HTML}{2F6FD6}\definecolor{fmB}{HTML}{EA6A10}\definecolor{fmR}{HTML}{E11D48}\definecolor{fmK}{HTML}{16A34A}\definecolor{fmV}{HTML}{7C3AED}

% lletres
\newcommand{\comic}{\fontfamily{ComicNeue-TLF}\selectfont}
\newcommand{\heavy}{\fontseries{eb}\selectfont}
\newcommand{\semi}{\fontseries{sb}\selectfont}
\DeclareUnicodeCharacter{2212}{-}

\setlength{\parindent}{0pt}
\setlength{\parskip}{7pt}
\setlength{\lineskip}{4pt}\setlength{\lineskiplimit}{4pt}
\setlength{\emergencystretch}{3em}
\setlist{itemsep=3pt,topsep=3pt,leftmargin=1.4em}
\setcounter{secnumdepth}{1}
\renewcommand{\thesection}{3.\arabic{section}}
\titleformat{\section}[hang]{\color{fmnavy}\fontsize{23}{27}\heavy\selectfont}{\tcbox[on line,colback=fmblue,colframe=fmblue,arc=3mm,boxsep=2pt,left=6pt,right=6pt,top=3pt,bottom=3pt]{\color{white}\thesection}}{0.45em}{}[\vspace{5pt}{\color{fmblue}\titlerule[2pt]}]
\titlespacing*{\section}{0pt}{4pt}{6pt}
\titleformat{\subsection}{\color{fmblue}\Large\semi}{}{0pt}{\faCaretRight\ }
\titlespacing*{\subsection}{0pt}{8pt}{2pt}

% frase de cada apartat, «a mà»
\newcommand{\frase}[1]{\par\vspace{-2pt}{\comic\Large\color{fmbluel}\faCommentDots[regular]\ \ #1\par}\vspace{2pt}}
\newcommand{\seccio}[3]{\Needspace{16\baselineskip}\section{#2}\hypertarget{#1}{}\frase{#3}}
% notes a mà dins dels exemples
\newcommand{\nota}[1]{{\comic\small\color{fmex!85!black}#1}}
\newcommand{\notab}[1]{{\comic\small\color{fmblue}#1}}
% bé / malament
\newcommand{\be}[1]{{\color{fmok}\faCheckCircle}\ \ensuremath{#1}}
\newcommand{\mal}[1]{{\color{fmR}\faTimesCircle}\ {\color{fmR}\ensuremath{#1}}}
% pas numerat
\newcommand{\pas}[1]{\tikz[baseline=(p.base)]\node[circle,fill=fmblue,text=white,inner sep=1.2pt,minimum size=1.25em,font=\small\heavy](p){#1};}
% etiqueta de color
\newcommand{\chip}[2]{\tcbox[on line,colback=#1!14,colframe=#1!60,arc=2mm,boxsep=1pt,left=5pt,right=5pt,top=1pt,bottom=1pt]{#2}}
\newcommand{\fmkey}[1]{\colorbox{fmyellow!55}{\strut #1}}
\newcommand{\cA}[1]{{\color{fmA}#1}}\newcommand{\cB}[1]{{\color{fmB}#1}}\newcommand{\cK}[1]{{\color{fmK}#1}}\newcommand{\cR}[1]{{\color{fmR}#1}}\newcommand{\cV}[1]{{\color{fmV}#1}}

% caixes
\tcbset{fmbase/.style={enhanced,arc=3mm,boxrule=1.2pt,left=4mm,right=4mm,top=2mm,bottom=2.5mm,before skip=8pt,after skip=8pt,fonttitle=\heavy\large,before upper=\raggedright}}
\newtcolorbox{idea}[1][Idea clau]{fmbase,colback=fmdefbg,colframe=fmdef,coltitle=white,title={\faKey\ \ #1}}
\newtcolorbox{exemple}[1][Exemple]{fmbase,colback=fmexbg,colframe=fmex,colbacktitle=fmexbg,coltitle=fmex,title={\faPencil*\ \ #1},titlerule=0.5pt,titlerule style={fmex!40,dashed}}
\newtcolorbox{compte}[1][Compte!]{fmbase,colback=fmsf!6,colframe=fmsf,coltitle=white,fonttitle=\comic\bfseries\Large,title={\faExclamationTriangle\ \ #1}}
\newtcolorbox{truc}[1][Truc]{fmbase,colback=fmyellow!20,colframe=fmgold,colbacktitle=fmyellow,coltitle=fmnavy,fonttitle=\comic\bfseries\Large,title={\faLightbulb[regular]\ \ #1}}

% dibuix de la targeta de la unitat, en blanc
\newcommand{\fmcard}{\begin{tikzpicture}[x=0.062cm,y=-0.062cm,white,line width=3pt,line cap=round]
  \begin{scope}[transparency group,opacity=0.15]
  \node[anchor=base west,inner sep=0pt,font=\fontsize{30}{30}\selectfont] at (14,44) {$x^2-5x+6>0$};
  \draw[line width=2pt] (120,95) -- (232,95);
  \draw[line width=3.2pt,domain=-40:40,samples=40] plot ({176+\x},{95-0.04*(\x*\x-400)});
  \fill (156,95) circle (5); \fill (196,95) circle (5);
  \end{scope}
\end{tikzpicture}}
\newcommand{\fmbrand}{\raisebox{-0.25\height}{\includegraphics[height=0.5cm]{logo-fentmates.pdf}}}

% regla de Ruffini: #1 = a, #2 = coeficients, #3 = fila del mig (sense el primer), #4 = fila de baix, #5 = nombre de coeficients
\newcommand{\ruffini}[5]{\begin{tikzpicture}[x=1.25cm,y=0.85cm,every node/.style={font=\large}]
  \foreach \v [count=\i] in {#2} \node (t\i) at (\i,0) {$\v$};
  \foreach \v [count=\i] in {#3} {\pgfmathtruncatemacro\j{\i+1}\node (m\j) at (\j,-1) {$\v$};}
  \foreach \v [count=\i] in {#4} \node (b\i) at (\i,-2.1) {$\v$};
  \node[text=fmB] at (0,-1) {$#1$};
  \draw[line width=1.1pt,fmnavy] (0.45,0.5)--(0.45,-1.5) (0.45,-1.5)--(#5+0.5,-1.5);
  \node[draw=fmR,line width=1.3pt,rounded corners=2mm,minimum size=0.75cm] at (#5,-2.1) {};
  \draw[-{Latex[length=2.2mm]},fmK,line width=1.1pt] (t1.south) -- (b1.north);
  \pgfmathtruncatemacro\last{#5-1}
  \foreach \i in {1,...,\last} {\pgfmathtruncatemacro\j{\i+1}
    \draw[-{Latex[length=2.2mm]},fmB,line width=1.1pt,shorten >=1pt] (b\i.north east) -- (m\j.south west);
    \draw[-{Latex[length=2.2mm]},fmA!70,line width=1pt,dashed,shorten >=2pt] ([xshift=8pt]t\j.south) -- ([xshift=8pt]b\j.north);}
\end{tikzpicture}}

% capçalera i peu
\pagestyle{fancy}\fancyhf{}
\fancyhead[L]{\small\fmbrand}
\fancyhead[R]{\small\color{fmmuted}4t ESO · Unitat 3: Equacions i inequacions}
\fancyfoot[L]{\scriptsize\color{fmmuted}\href{https://creativecommons.org/licenses/by-nc-sa/4.0/deed.ca}{\ccbyncsa}\ \ FentMates · fentmates.org · CC BY-NC-SA 4.0}
\fancyfoot[R]{\small\color{fmmuted}\thepage}
\renewcommand{\headrulewidth}{0.4pt}\renewcommand{\headrule}{\color{fmmuted!40}\hrule height \headrulewidth}
\fancypagestyle{primera}{\fancyhf{}\renewcommand{\headrulewidth}{0pt}
  \fancyfoot[L]{\scriptsize\color{fmmuted}\href{https://creativecommons.org/licenses/by-nc-sa/4.0/deed.ca}{\ccbyncsa}\ \ FentMates · fentmates.org · CC BY-NC-SA 4.0}
  \fancyfoot[R]{\small\color{fmmuted}\thepage}}

% targeta de l'índex
\newcommand{\itile}[4]{\begin{tcolorbox}[enhanced,colback=white,colframe=fmblue!35,arc=3mm,boxrule=1pt,boxsep=0pt,left=6pt,right=6pt,top=5pt,bottom=5pt,before skip=0pt,after skip=7pt]
  \hyperlink{#1}{\begin{minipage}[c]{0.9cm}\centering{\color{fmblue}\LARGE #3}\end{minipage}\hspace{6pt}%
  \begin{minipage}[c]{0.78\linewidth}{\color{fmblue}\heavy #2}\\[-2pt]{\color{fmnavy}\semi\large #4}\end{minipage}}\end{tcolorbox}}

% taula amb vores arrodonides: el contingut es retalla amb la forma (sense @{} a les vores ni entre columnes)
\newsavebox\tabclipbox
\newcommand{\taulaclip}[3][3mm]{\sbox\tabclipbox{#3}\begin{tikzpicture}[baseline=(tcn.base)]
  \begin{scope}\clip[rounded corners=#1] (0,-\dp\tabclipbox) rectangle (\wd\tabclipbox,\ht\tabclipbox);
  \node[inner sep=0pt,outer sep=0pt,anchor=base west] (tcn) at (0,0) {\usebox\tabclipbox};\end{scope}
  \draw[#2,rounded corners=#1] (0,-\dp\tabclipbox) rectangle (\wd\tabclipbox,\ht\tabclipbox);
\end{tikzpicture}}
\newcommand{\taularod}[3][fmblue]{\taulaclip{#1,line width=1.2pt}{\renewcommand{\arraystretch}{1.5}\begin{tabular}{#2}#3\end{tabular}}}
% balança
\newcommand{\balanca}[3]{\begin{tikzpicture}[scale=#3,transform shape,every node/.style={anchor=center}]
  \fill[fmnavy] (-0.45,-1.9) rectangle (0.45,-1.75);
  \fill[fmnavy!80] (-0.12,-1.75) rectangle (0.12,0);
  \fill[fmgold] (0,0.05) circle (0.13);
  \draw[fmnavy,line width=2.4pt,line cap=round] (-2,0) -- (2,0);
  \foreach \s in {-1,1} {
    \draw[fmmuted,line width=0.8pt] (\s*2,0) -- (\s*2-0.75,-1.05) (\s*2,0) -- (\s*2+0.75,-1.05);
    \fill[fmyellow!70,draw=fmgold,line width=1pt] (\s*2-0.95,-1.05) arc[start angle=180,end angle=360,x radius=0.95,y radius=0.32] -- cycle;}
  \node[font=\Large,fill=fmA!15,draw=fmA,rounded corners=1.5mm,inner sep=3pt] at (-2,-0.68) {$#1$};
  \node[font=\Large,fill=fmB!15,draw=fmB,rounded corners=1.5mm,inner sep=3pt] at (2,-0.68) {$#2$};
\end{tikzpicture}}
% eixos amb quadrícula (dins d'un tikzpicture): xmin, xmax, ymin, ymax
\newcommand{\eixos}[4]{\draw[fmmuted!22,line width=0.4pt] (#1,#3) grid (#2,#4);
  \draw[-{Latex[length=2mm]},fmmuted,line width=0.8pt] (#1-0.2,0)--(#2+0.4,0) node[right,font=\small]{$x$};
  \draw[-{Latex[length=2mm]},fmmuted,line width=0.8pt] (0,#3-0.2)--(0,#4+0.4) node[above,font=\small]{$y$};}
% taula de traduccions (diccionari) amb les vores arrodonides
\newcommand{\dic}[2]{#1 & {\large$#2$}\\}
\newcommand{\dicratio}{0.62}
\newcommand{\dichead}{\begin{tabular}{@{}>{\raggedright}p{\dimexpr\dicratio\linewidth-12pt}>{\centering\arraybackslash}p{\dimexpr\linewidth-\dicratio\linewidth-12pt}@{}}\faComment[regular]\ Ho diem així\dots & \dots\ així\end{tabular}}
\newenvironment{dicbox}{\begin{tcolorbox}[enhanced,arc=3mm,boxrule=1.2pt,colframe=fmblue,colback=white,coltitle=white,fonttitle=\heavy,
  title={\dichead},left=0pt,right=0pt,top=0pt,bottom=1.5mm,boxsep=0pt,lefttitle=6pt,righttitle=6pt,before skip=0pt,after skip=0pt]\rowcolors{1}{white}{fmsurface}%
  \begin{tabular}{@{\hspace{6pt}}>{\raggedright}p{\dimexpr\dicratio\linewidth-12pt}>{\centering\arraybackslash}p{\dimexpr\linewidth-\dicratio\linewidth-12pt}@{\hspace{6pt}}}}
  {\end{tabular}\end{tcolorbox}}
% resum
\newcommand{\rtile}[3]{\begin{tcolorbox}[enhanced,colback=white,colframe=fmblue!45,arc=3mm,boxrule=1pt,left=3mm,right=3mm,top=2.5mm,bottom=2.5mm,fontupper=\large,
  title={\tcbox[on line,colback=fmblue,colframe=fmblue,arc=2mm,boxsep=0pt,left=4pt,right=4pt,top=1pt,bottom=1pt]{\color{white}\small\heavy #1}\ \ #2},coltitle=fmnavy,colbacktitle=fmdefbg,fonttitle=\semi,toptitle=0.5mm,bottomtitle=0.5mm]#3\end{tcolorbox}}
% portada: #1 subtítol a mà
\newcommand{\portada}[1]{\thispagestyle{primera}
\begin{tikzpicture}[remember picture,overlay]
  \fill[fmblue] (current page.north west) rectangle ([yshift=-9.4cm]current page.north east);
  \begin{scope}
    \clip (current page.north west) rectangle ([yshift=-9.4cm]current page.north east);
    \node[anchor=north west,inner sep=0pt] at ([xshift=6.6cm,yshift=-1.3cm]current page.north west) {\fmcard};
  \end{scope}
  \node[anchor=north west,inner sep=0pt] at ([xshift=1.6cm,yshift=-1.1cm]current page.north west) {\tcbox[colback=white,colframe=white,arc=2.5mm,boxsep=0pt,left=7pt,right=7pt,top=5pt,bottom=5pt]{\includegraphics[height=0.75cm]{logo-fentmates.pdf}}};
\end{tikzpicture}
\vspace*{0.55cm}
\tcbox[colback=white!15!fmblue,colframe=white!15!fmblue,arc=2.5mm,boxsep=0pt,left=6pt,right=6pt,top=3pt,bottom=3pt]{\color{white}\small\bfseries Apunts de teoria · 4t ESO · Matemàtiques B}

\vspace{2pt}
{\raggedright\color{white}\fontsize{30}{36}\heavy\selectfont{\color{fmyellow}Unitat 3}\\[2pt]Equacions i inequacions\par}
\vspace{8pt}
{\color{white}\comic\Large #1\par}
\vspace*{\dimexpr 7.4cm-\pagetotal\relax}
}
\newcommand{\llegenda}{\par\vspace{16pt}
{\color{fmnavy}\heavy\Large\faPalette\ \ Com es llegeixen aquests apunts}\par\smallskip
\begin{tcbraster}[raster columns=4,raster equal height,raster column skip=8pt,raster before skip=12pt]
\tcbox[colback=fmdefbg,colframe=fmdef,arc=3mm,boxrule=1.2pt,halign=center,width=\linewidth]{{\color{fmdef}\Large\faKey}\\[2pt]{\color{fmdef}\heavy Idea clau}\\[-1pt]{\small el que cal saber}}
\tcbox[colback=fmexbg,colframe=fmex,arc=3mm,boxrule=1.2pt,halign=center,width=\linewidth]{{\color{fmex}\Large\faPencil*}\\[2pt]{\color{fmex}\heavy Exemple}\\[-1pt]{\small pas a pas}}
\tcbox[colback=fmsf!6,colframe=fmsf,arc=3mm,boxrule=1.2pt,halign=center,width=\linewidth]{{\color{fmsf}\Large\faExclamationTriangle}\\[2pt]{\color{fmsf}\heavy Compte!}\\[-1pt]{\small errors típics}}
\tcbox[colback=fmyellow!20,colframe=fmgold,arc=3mm,boxrule=1.2pt,halign=center,width=\linewidth]{{\color{fmgold!80!black}\Large\faLightbulb[regular]}\\[2pt]{\color{fmnavy}\heavy Truc}\\[-1pt]{\small per a recordar-ho}}
\end{tcbraster}}
\newcommand{\resumtitol}{\clearpage
\hypertarget{resum}{}
{\color{fmnavy}\fontsize{23}{27}\heavy\selectfont\faClipboardList\ \ En resum}\par\vspace{5pt}{\color{fmblue}\titlerule[2pt]}
\frase{Tota la unitat en una pàgina.}}
% els 5 passos per a resoldre problemes: #2 = text del pas 2
\newcommand{\passosproblema}[2]{\begin{center}
\begin{tikzpicture}[s/.style={draw=fmblue,fill=fmdefbg,rounded corners=3mm,line width=1.2pt,minimum width=2.75cm,minimum height=2.45cm,align=center,font=\small},
  ic/.style={font=\LARGE,text=fmblue}]
  \foreach \i/\ic/\t in {1/\faBookOpen/{\textbf{Llegim}\\i entenem},2/\faQuestionCircle/{#1},3/\faEquals/{#2},4/\faCogs/{\textbf{Resolem}},5/\faCheckDouble/{\textbf{Comprovem}\\i responem}} {
    \node[s] (s\i) at (3.3*\i,0) {};
    \node[ic] at ([yshift=-0.5cm]s\i.north) {\ic};
    \node[align=center,font=\small,anchor=north] at ([yshift=-0.95cm]s\i.north) {\t};
    \node[circle,fill=fmblue,text=white,font=\small\heavy,inner sep=1.5pt,minimum size=1.3em] at (s\i.north west) {\i};}
  \foreach \i [evaluate=\i as \j using int(\i+1)] in {1,...,4} \draw[-{Latex[length=3mm]},fmmuted,line width=1.5pt] (s\i) -- (s\j);
\end{tikzpicture}
\end{center}}

\newcommand{\lo}{\mathopen{]}}\newcommand{\ro}{\mathclose{[}}
% punts ple i buit, nombres dels eixos i taula de valors
\newcommand{\pple}[2][fmblue]{\fill[#1] #2 circle (3.6pt);}
\newcommand{\nums}[4]{\foreach \k in {#1,...,#2} {\ifnum\k=0\else\node[below,font=\scriptsize,text=fmmuted] at (\k,0) {$\k$};\fi}
  \foreach \k in {#3,...,#4} {\ifnum\k=0\else\node[left,font=\scriptsize,text=fmmuted] at (0,\k) {$\k$};\fi}}
\newcommand{\tv}[3][fmblue]{\taulaclip{#1,line width=1.2pt}{\renewcommand{\arraystretch}{1.45}%
  \begin{tabular}{>{\columncolor{#1}\color{white}\heavy}c*{12}{w{c}{0.8cm}}}#2\\\rowcolor{#1!8}\cellcolor{#1}#3\end{tabular}}}
% etiqueta d'una recta o corba
\newcommand{\etq}[3]{\node[fill=white,inner sep=1.5pt,text=#1,font=\small] at #2 {#3};}

\newcommand{\passos}[1]{\begin{tcolorbox}[fmbase,colback=fmdefbg,colframe=fmdef,coltitle=white,title={\faListOl\ \ Com es resolen}]#1\end{tcolorbox}}

\begin{document}
\portada{Equacions de tota mena i desigualtats: trobem quins valors les compleixen.}

{\color{fmnavy}\heavy\Large\faMapSigns\ \ Què aprendràs?}\par\smallskip
\begin{minipage}[t]{0.485\linewidth}
\itile{sec1}{3.1}{\faSuperscript}{Equacions de segon grau}
\itile{sec2}{3.2}{\faIcon{exchange-alt}}{Biquadrades}
\itile{sec3}{3.3}{\faCubes}{Grau major que 2}
\itile{sec4}{3.4}{\faPuzzlePiece}{Equacions factoritzades}
\itile{sec5}{3.5}{\faDivide}{Amb fraccions algebraiques}
\end{minipage}\hfill
\begin{minipage}[t]{0.485\linewidth}
\itile{sec6}{3.6}{\faIcon{square-root-alt}}{Amb radicals}
\itile{sec7}{3.7}{\faRocket}{Exponencials}
\itile{sec8}{3.8}{\faIcon{calculator}}{Logarítmiques}
\itile{sec9}{3.9}{\faAngleRight}{Inequacions de primer grau}
\itile{sec10}{3.10}{\faChartArea}{Inequacions de grau 2 o més}
\end{minipage}
\llegenda

% ============================================================ 3.1
\clearpage
\seccio{sec1}{Equacions de segon grau}{La fórmula, el discriminant i les incompletes.}

\vfill
\begin{minipage}[c]{0.5\linewidth}
\begin{tcolorbox}[enhanced,colback=fmyellow!18,colframe=fmgold,arc=4mm,boxrule=1.6pt,before upper=\centering,top=3mm,bottom=3mm]
{\large$ax^2+bx+c=0$}\par\medskip
{\fontsize{20}{26}\selectfont$x=\dfrac{-b\pm\sqrt{\cR{b^2-4ac}}}{2a}$}
\end{tcolorbox}
\end{minipage}\hfill
\begin{minipage}[c]{0.47\linewidth}
\begin{idea}[El discriminant $\Delta=b^2-4ac$]
\chip{fmK}{$\Delta>0$} dues solucions\\[3pt]
\chip{fmgold}{$\Delta=0$} una doble, $x=-\frac{b}{2a}$\\[3pt]
\chip{fmR}{$\Delta<0$} cap solució real
\end{idea}
\end{minipage}

\vfill
\begin{tcbraster}[raster columns=2,raster equal height,raster column skip=10pt]
\begin{tcolorbox}[fmbase,colback=fmA!8,colframe=fmA,coltitle=white,title={Sense terme independent}]
$ax^2+bx=0\ \to\ x(ax+b)=0$\par\smallskip
$x=0$, \ $x=-\dfrac ba$
\end{tcolorbox}
\begin{tcolorbox}[fmbase,colback=fmB!8,colframe=fmB,coltitle=white,title={Sense terme en $x$}]
$ax^2+c=0\ \to\ x^2=-\dfrac ca$\par\smallskip
$x=\pm\sqrt{-\frac ca}$ \ (si és $\ge0$)
\end{tcolorbox}
\end{tcbraster}

\vfill
\begin{compte}[No dividisques entre $x$]
En $x^2-5x=0$, si divideixes entre $x$ perds la solució $x=0$. \textbf{Trau factor comú}.
\end{compte}

\vfill
\begin{exemple}[Exemples]
$2x^2-x-3=0$: \ $\Delta=1+24=25$, \ $x=\dfrac{1\pm5}{4}$ \ $\to$ \ $\boxed{x=\tfrac32,\ x=-1}$\\[5pt]
$3x^2-12=0\ \to\ x^2=4\ \to\ \boxed{x=\pm2}$ \qquad $5x^2+10x=0\ \to\ 5x(x+2)=0\ \to\ \boxed{x=0,\ x=-2}$\\[5pt]
$(x-2)^2=2x-1\ \to\ x^2-6x+5=0\ \to\ \boxed{x=5,\ x=1}$ \ \nota{primer, la forma $ax^2+bx+c=0$}\\[5pt]
$4x^2-12x+9=0$: \ $\Delta=0$, \ $\boxed{x=\tfrac32}$ doble \qquad $x^2+x+5=0$: \ $\Delta=1-20<0$, \ \textbf{cap solució}
\end{exemple}

\vfill
\begin{exemple}[Exemple amb parèntesis: \ $(2x-1)^2-(x+2)(x-2)=4x+1$]
\pas{1} Desfem els parèntesis: \ $4x^2-4x+1-(x^2-4)=4x+1$ \ \nota{el $-$ canvia tots els signes}\\[4pt]
\pas{2} Tot a un membre: \ $3x^2-8x+4=0$ \qquad \pas{3} $x=\dfrac{8\pm\sqrt{64-48}}{6}=\dfrac{8\pm4}{6}$ \ $\to$ \ \boxed{x=2,\ x=\tfrac23}
\end{exemple}

\vfill
\begin{exemple}[Exemple amb denominadors: \ $\dfrac{x(x-1)}{2}-\dfrac{(x+1)^2}{3}=-2$]
\pas{1} m.c.m. $=6$: \ $3x(x-1)-2(x+1)^2=-12$\\[4pt]
\pas{2} $3x^2-3x-2x^2-4x-2=-12\ \to\ x^2-7x+10=0$ \qquad \pas{3} $x=\dfrac{7\pm3}{2}$ \ $\to$ \ \boxed{x=5,\ x=2}
\end{exemple}

\vfill
% ============================================================ 3.2
\clearpage
\seccio{sec2}{Equacions biquadrades}{Només graus 4, 2 i 0: canvi $t=x^2$.}

\vfill
\begin{minipage}[c]{0.5\linewidth}
\passos{\pas{1} Canvi \ $t=x^2$ \ (i $x^4=t^2$): \ $at^2+bt+c=0$.\\[3pt]
\pas{2} Resolem en $t$.\\[3pt]
\pas{3} Desfem el canvi: \ $t>0\to x=\pm\sqrt t$; \ $t=0\to x=0$; \ $t<0$: cap.}
\end{minipage}\hfill
\begin{minipage}[c]{0.47\linewidth}
\begin{truc}[Quantes solucions?]
Una biquadrada pot tindre \textbf{4, 3, 2, 1 o cap} solució.
\end{truc}
\end{minipage}

\vfill
\begin{exemple}[Exemple: \ $x^4-5x^2+4=0$]
$t^2-5t+4=0\ \to\ t=\dfrac{5\pm3}{2}$: \ $t=4,\ t=1$ \qquad $x^2=4\to x=\pm2$; \ $x^2=1\to x=\pm1$ \qquad \boxed{\text{quatre solucions}}
\end{exemple}

\vfill
\begin{exemple}[Exemple: \ $x^4+3x^2-4=0$]
$t^2+3t-4=0\ \to\ t=1,\ t=-4$ \qquad $x^2=1\to\boxed{x=\pm1}$; \ $x^2=-4$ \mal{\text{no té solució}}
\end{exemple}

\vfill
\begin{exemple}[Exemples]
$x^4-13x^2+36=0$: \ $t=9,\ t=4$ \ $\to$ \ $x=\pm3,\ \pm2$ \qquad $x^4-16=0$: \ $t=4$ \ $\to$ \ $x=\pm2$\\[4pt]
$x^4+5x^2+6=0$: \ $t=-2,\ -3$: \ \textbf{cap} \qquad $x^4-3x^2=0$: \ $x^2(x^2-3)=0$ \ $\to$ \ $x=0,\ \pm\sqrt3$
\end{exemple}

\vfill
\begin{exemple}[Exemple amb parèntesis: \ $(x^2-2)(x^2+1)=10$]
$x^4+x^2-2x^2-2=10\ \to\ x^4-x^2-12=0$ \qquad $t^2-t-12=0\ \to\ t=4,\ t=-3$ \qquad $x^2=4\to\boxed{x=\pm2}$; \ $x^2=-3$ \mal{\text{no}}
\end{exemple}

\vfill
\begin{exemple}[Exemple amb denominadors: \ $\dfrac{x^4}{4}-\dfrac{5x^2}{2}+4=0$]
Per 4: \ $x^4-10x^2+16=0$ \qquad $t=\dfrac{10\pm6}{2}=8,\ 2$ \qquad $x^2=8\to x=\pm2\sqrt2$; \ $x^2=2\to x=\pm\sqrt2$ \qquad \boxed{\text{quatre solucions}}
\end{exemple}

\vfill
\begin{minipage}[t]{0.48\linewidth}
\begin{exemple}[Exemple: \ $x^2-\dfrac{36}{x^2}=5$]
Per $x^2$: \ $x^4-5x^2-36=0$\\[3pt]
$t=9,\ t=-4$ \ $\to$ \ \boxed{x=\pm3}
\end{exemple}
\end{minipage}\hfill
\begin{minipage}[t]{0.49\linewidth}
\begin{truc}[El mateix canvi, amb $x^3$]
$x^6-9x^3+8=0$: \ $t=x^3$\\[3pt]
$t^2-9t+8=0\ \to\ t=8,\ t=1$ \ $\to$ \ $x=2,\ x=1$
\end{truc}
\end{minipage}

\vfill
% ============================================================ 3.3
\clearpage
\seccio{sec3}{Equacions de grau major que 2}{Factoritzem amb Ruffini fins a grau 2.}

\vfill
\begin{minipage}[c]{0.5\linewidth}
\passos{\pas{1} Tot a un membre i, si es pot, factor comú.\\[3pt]
\pas{2} Ruffini amb els \textbf{divisors del terme independent}, fins a grau 2.\\[3pt]
\pas{3} Resolem la de segon grau. Les solucions són totes les arrels.}
\end{minipage}\hfill
\begin{minipage}[c]{0.47\linewidth}\centering
{\small$x^3-2x^2-5x+6=0$}\\[2pt]
\ruffini{1}{1,-2,-5,6}{1,-1,-6}{1,-1,-6,0}{4}
\end{minipage}

\vfill
\begin{exemple}[Exemple: \ $x^3-2x^2-5x+6=0$]
Candidats $\pm1,\pm2,\pm3,\pm6$. \ Amb $x=1$, residu 0: \ $(x-1)(x^2-x-6)=0$\\[3pt]
$x^2-x-6=0\ \to\ x=\dfrac{1\pm5}{2}$ \qquad Solucions: \ \boxed{x=1,\ x=3,\ x=-2}
\end{exemple}

\vfill
\begin{exemple}[Exemple: \ $x^4=x^3+4x^2-4x$]
Tot a un membre i factor comú: \ $x(x^3-x^2-4x+4)=0$ \qquad Ruffini amb $x=1$: \ $x(x-1)(x^2-4)=0$\\[3pt]
\boxed{x=0,\ x=1,\ x=\pm2}
\end{exemple}

\vfill
\begin{exemple}[Exemples]
$x^3-7x+6=0$: \ $(x-1)(x-2)(x+3)=0$ \qquad $x^3+x^2-5x+3=0$: \ $(x-1)^2(x+3)=0$ \ \nota{$x=1$ doble}\\[4pt]
$2x^3-3x^2-3x+2=0$: \ $(x-2)(2x^2+x-1)=0$ \ $\to$ \ $x=2,\ \frac12,\ -1$\\[4pt]
$x^3+2x^2+2x+1=0$: \ $(x+1)(x^2+x+1)=0$ \ $\to$ \ només $x=-1$ \ \nota{$\Delta=1-4<0$}
\end{exemple}

\vfill
\begin{exemple}[Exemple: \ $x^4-2x^3-3x^2+4x+4=0$ \ (Ruffini dues vegades)]
\begin{minipage}[c]{0.4\linewidth}\centering
\scalebox{0.62}{\ruffini{-1}{1,-2,-3,4,4}{-1,3,0,-4}{1,-3,0,4,0}{5}}
\end{minipage}\hfill
\begin{minipage}[c]{0.57\linewidth}
$x=-1$: \ $(x+1)(x^3-3x^2+4)=0$\\[3pt]
Ruffini amb $x=2$: \ $x^3-3x^2+4=(x-2)(x^2-x-2)$\\[3pt]
$x^2-x-2=0\ \to\ x=2,\ x=-1$\\[3pt]
$(x+1)^2(x-2)^2=0$: \ \boxed{x=-1,\ x=2} \ \nota{totes dues dobles}
\end{minipage}
\end{exemple}

\vfill
\begin{compte}[No dividisques entre $x$]
$x^3=4x$: \ si dividim entre $x$ perdem el 0. \ Tot a un membre: \ $x(x^2-4)=0$ \ $\to$ \ $x=0,\ x=\pm2$
\end{compte}

\vfill
% ============================================================ 3.4
\clearpage
\seccio{sec4}{Equacions factoritzades}{Si un producte val 0, algun factor val 0.}

\vfill
\begin{tcolorbox}[enhanced,colback=fmyellow!18,colframe=fmgold,arc=4mm,boxrule=1.6pt,before upper=\centering,top=3mm,bottom=3mm]
{\Large$A\cdot B\cdot C=0\iff A=0\ \ \text{o}\ \ B=0\ \ \text{o}\ \ C=0$}\par\smallskip
{\small no cal multiplicar: igualem cada factor a zero}
\end{tcolorbox}

\vfill
\begin{compte}[Només si l'altre membre és 0]
En $(x-1)(x+2)=4$ \textbf{no} podem dir $x-1=4$. Cal multiplicar i resoldre $x^2+x-6=0$.
\end{compte}

\vfill
\begin{exemple}[Exemple: \ $(x-2)(x+3)(2x-1)=0$]
$x-2=0\to x=2$; \quad $x+3=0\to x=-3$; \quad $2x-1=0\to x=\frac12$
\end{exemple}

\vfill
\begin{exemple}[Exemple: \ $(x^2-9)(x^2+x-2)=0$]
$x^2-9=0\to x=\pm3$ \qquad $x^2+x-2=0\to x=1,\ x=-2$ \qquad Solucions: \ \boxed{-3,\ -2,\ 1,\ 3}
\end{exemple}

\vfill
\begin{exemple}[Exemples]
$x^2(x-4)(x^2+1)=0$: \ $x=0$ (doble), \ $x=4$; \ $x^2+1=0$ no té solució\\[4pt]
$(2x+1)^2(x^2+4)=0$: \ $x=-\frac12$ (doble) \qquad $(x+3)(x-1)=5$: \ $x^2+2x-8=0\to x=2,\ x=-4$\\[4pt]
$(x^2-4x+4)(x^3-x)=0$: \ $(x-2)^2\,x\,(x+1)(x-1)=0$ \ $\to$ \ $x=2$ (doble), \ $0,\ \pm1$
\end{exemple}

\vfill
\begin{exemple}[Exemple: \ $(x^2-1)^2-(x^2-1)(x+1)=0$ \ (factor comú amb parèntesis)]
Traiem factor comú $(x^2-1)$: \ $(x^2-1)\big[(x^2-1)-(x+1)\big]=0\ \to\ (x^2-1)(x^2-x-2)=0$\\[4pt]
$x^2-1=0\to x=\pm1$ \qquad $x^2-x-2=0\to x=2,\ x=-1$ \qquad Solucions: \ \boxed{x=-1\text{ (doble)},\ 1,\ 2}
\end{exemple}

\vfill
\begin{minipage}[t]{0.48\linewidth}
\begin{compte}[Tampoc simplifiques]
$x^2(x-3)=4x^2$: \ si dividim entre $x^2$ perdem $x=0$.\\[3pt]
$x^2(x-3)-4x^2=0\ \to\ x^2(x-7)=0$\\[2pt]
\boxed{x=0,\ x=7}
\end{compte}
\end{minipage}\hfill
\begin{minipage}[t]{0.49\linewidth}
\begin{exemple}[Exemple: \ solucions $2$, $-1$ i $\frac13$]
Un factor per solució:\\[3pt]
$(x-2)(x+1)(3x-1)=0$\\[2pt]
\notab{$x=\frac13$ dona el factor $3x-1$}
\end{exemple}
\end{minipage}

\vfill
% ============================================================ 3.5
\clearpage
\seccio{sec5}{Equacions amb fraccions algebraiques}{Multipliquem pel m.c.m. i comprovem.}

\vfill
\begin{minipage}[c]{0.55\linewidth}
\passos{\pas{1} Factoritzem els denominadors i en fem el \textbf{m.c.m.}\\[3pt]
\pas{2} Multipliquem els dos membres pel m.c.m.\\[3pt]
\pas{3} Resolem l'equació polinòmica.\\[3pt]
\pas{4} \textbf{Comprovem}: descartem les que anul·len un denominador.}
\end{minipage}\hfill
\begin{minipage}[c]{0.42\linewidth}
\begin{compte}[Pot no tindre solució]
$\dfrac{x}{x-2}-\dfrac{2}{x-2}=3\ \to\ x=2$, però anul·la el denominador: \textbf{no té solució}.
\end{compte}
\end{minipage}

\vfill
\begin{exemple}[Exemple: \ $\dfrac1x+\dfrac{1}{x+2}=\dfrac34$]
m.c.m. $=4x(x+2)$: \ $4(x+2)+4x=3x(x+2)\ \to\ 3x^2-2x-8=0\ \to\ x=\dfrac{2\pm10}{6}$\\[3pt]
\boxed{x=2,\ x=-\tfrac43} \ \nota{cap anul·la els denominadors}
\end{exemple}

\vfill
\begin{exemple}[Exemple: \ $\dfrac{x}{x-1}+\dfrac{2}{x+1}=\dfrac{2}{x^2-1}$]
m.c.m. $=(x+1)(x-1)$: \ $x(x+1)+2(x-1)=2\ \to\ x^2+3x-4=0\ \to\ x=1,\ x=-4$\\[3pt]
$x=1$ anul·la $x-1$: es descarta. \ Solució: \boxed{x=-4}
\end{exemple}

\vfill
\begin{exemple}[Exemples]
$\dfrac3x+2=5\ \to\ x=1$ \qquad $\dfrac{x}{x+1}=\dfrac2x\ \to\ x^2-2x-2=0\ \to\ x=1\pm\sqrt3$\\[5pt]
$\dfrac{6}{x^2-4}-\dfrac{1}{x-2}=\dfrac{1}{x+2}$: \ m.c.m. $(x+2)(x-2)$: \ $6-(x+2)=x-2\ \to\ x=3$
\end{exemple}

\vfill
\begin{exemple}[Exemple: \ $\dfrac{x+1}{x-1}-\dfrac{x-1}{x+1}=\dfrac83$]
m.c.m. $=3(x+1)(x-1)$: \ $3(x+1)^2-3(x-1)^2=8(x^2-1)$\\[3pt]
$3(x^2+2x+1)-3(x^2-2x+1)=8x^2-8\ \to\ 12x=8x^2-8\ \to\ 2x^2-3x-2=0$ \ $\to$ \ $x=\dfrac{3\pm5}{4}$: \ \boxed{x=2,\ x=-\tfrac12}
\end{exemple}

\vfill
\newpage
\vfill
\begin{exemple}[Exemple: \ $\dfrac2x+\dfrac{x}{x-3}=\dfrac{9}{x^2-3x}$]
\begin{minipage}[t]{0.6\linewidth}
Factoritzem: \ $x^2-3x=x(x-3)$, \ que ja és el m.c.m.\\[3pt]
$2(x-3)+x\cdot x=9\ \to\ x^2+2x-15=0\ \to\ x=3,\ x=-5$\\[3pt]
$x=3$ anul·la $x-3$ \ \mal{x=3} \qquad Solució: \boxed{x=-5}
\end{minipage}\hfill
\begin{minipage}[t]{0.36\linewidth}
\notab{Comprovem $x=-5$:}\\[2pt]
$-\frac25+\frac58=\frac{9}{40}$ \ i \ $\frac{9}{25+15}=\frac{9}{40}$ \ \be{}
\end{minipage}
\end{exemple}

\vfill
\begin{truc}[Abans del m.c.m., factoritza]
$x^2-3x=x(x-3)$, \ $x^2-1=(x+1)(x-1)$, \ $x^2-4x+4=(x-2)^2$\dots\ Així el m.c.m. és més senzill i veus de seguida quins valors anul·len els denominadors.
\end{truc}

\vfill
\subsection{Problemes que porten a fraccions}

\vfill
\begin{exemple}[Problema: \ un grup paga 120 € a parts iguals. Si foren 2 més, cadascú pagaria 5 € menys. Quants són?]
$x$ persones: \ $\dfrac{120}{x}-\dfrac{120}{x+2}=5\ \to\ 120(x+2)-120x=5x(x+2)\ \to\ x^2+2x-48=0\ \to\ x=6,\ x=-8$\\[3pt]
Una quantitat de persones no pot ser negativa: \ \boxed{6\text{ persones}} \ \nota{$20$ € cadascú; amb 8, $15$ €}
\end{exemple}

\vfill
\begin{exemple}[Problema: \ una aixeta omple un dipòsit en unes hores i una altra tarda 3 h més. Juntes, l'omplin en 2 h. Quant tarda cadascuna?]
En 1 hora: \ $\dfrac1x+\dfrac{1}{x+3}=\dfrac12$ \ \nota{la part del dipòsit que omplin}\\[4pt]
m.c.m. $2x(x+3)$: \ $2(x+3)+2x=x(x+3)\ \to\ x^2-x-6=0\ \to\ x=3,\ x=-2$ \qquad \boxed{3\text{ h i }6\text{ h}}
\end{exemple}

\vfill
\begin{exemple}[Problema: \ un ciclista fa 60 km. Si anara 5 km/h més ràpid, tardaria 1 hora menys. A quina velocitat va?]
Temps $=\dfrac{\text{distància}}{\text{velocitat}}$: \ $\dfrac{60}{x}-\dfrac{60}{x+5}=1\ \to\ 60(x+5)-60x=x(x+5)\ \to\ x^2+5x-300=0$\\[4pt]
$x=\dfrac{-5\pm35}{2}=15,\ -20$ \qquad \boxed{15\text{ km/h}} \ \nota{tarda 4 h; a 20 km/h, 3 h}
\end{exemple}

\vfill
% ============================================================ 3.6
\clearpage
\seccio{sec6}{Equacions amb radicals}{Aïllem l'arrel, elevem al quadrat i comprovem.}

\vfill
\begin{minipage}[c]{0.55\linewidth}
\passos{\pas{1} \textbf{Aïllem} el radical en un membre.\\[3pt]
\pas{2} \textbf{Elevem al quadrat} els dos membres.\\[3pt]
\pas{3} Si queda una arrel, repetim.\\[3pt]
\pas{4} Resolem i \textbf{comprovem} totes les solucions.}
\end{minipage}\hfill
\begin{minipage}[c]{0.42\linewidth}
\begin{compte}[Solucions falses]
En elevar al quadrat poden aparéixer solucions que no ho són: \ $2\neq-2$, però $2^2=(-2)^2$.
\end{compte}
\end{minipage}

\vfill
\begin{exemple}[Exemple: \ $x-\sqrt{x+1}=5$]
Aïllem: \ $x-5=\sqrt{x+1}$ \qquad Al quadrat: \ $x^2-10x+25=x+1\ \to\ x^2-11x+24=0\ \to\ x=8,\ x=3$\\[3pt]
Comprovem: \ $8-3=5$ \be{} \qquad $3-2=1\neq5$ \ \mal{x=3} \qquad Solució: \boxed{x=8}
\end{exemple}

\vfill
\begin{exemple}[Exemple: \ $\sqrt{2x+3}-\sqrt{x+1}=1$ \ \ (dos radicals)]
$\sqrt{2x+3}=1+\sqrt{x+1}\ \to\ 2x+3=1+2\sqrt{x+1}+x+1\ \to\ x+1=2\sqrt{x+1}$\\[3pt]
$x^2+2x+1=4x+4\ \to\ x^2-2x-3=0\ \to\ x=3,\ x=-1$ \qquad totes dues valen: \ \boxed{x=3,\ x=-1}
\end{exemple}

\vfill
\begin{exemple}[Exemples]
$\sqrt{3x+1}=4\ \to\ x=5$ \be{} \qquad $\sqrt{2x-1}=x-2\ \to\ x^2-6x+5=0$: \ $x=5$ \be{}, \ $x=1$ \mal{1\neq-1}
\end{exemple}

\vfill
\begin{truc}[El quadrat d'una suma]
$(1+\sqrt{x+1})^2=1+2\sqrt{x+1}+(x+1)$: \ no t'oblides del \textbf{doble producte}.
\end{truc}

\vfill
\begin{exemple}[Exemple: \ $\sqrt{x+5}+\sqrt x=5$ \ (deixem un radical a cada membre)]
$\sqrt{x+5}=5-\sqrt x\ \to\ x+5=25-10\sqrt x+x\ \to\ 10\sqrt x=20\ \to\ \sqrt x=2\ \to\ \boxed{x=4}$ \qquad Comprovem: \ $\sqrt9+\sqrt4=5$ \ \be{}
\end{exemple}

\vfill
\begin{minipage}[t]{0.48\linewidth}
\begin{exemple}[Problema: \ un nombre més la seua arrel quadrada fa 12]
$x+\sqrt x=12\ \to\ \sqrt x=12-x$\\[3pt]
$x=144-24x+x^2\ \to\ x^2-25x+144=0$\\[3pt]
$x=16$: \ $16+4=20$ \ \mal{\text{no}}\\[2pt] $x=9$: \ $9+3=12$ \ \be{}\\[3pt]
El nombre és \boxed{9}
\end{exemple}
\end{minipage}\hfill
\begin{minipage}[t]{0.49\linewidth}
\begin{compte}[Una arrel no és mai negativa]
$\sqrt{x-1}+3=1\ \to\ \sqrt{x-1}=-2$\\[3pt]
\textbf{No té solució}: no cal ni elevar al quadrat.
\end{compte}
\end{minipage}

\vfill
% ============================================================ 3.7
\clearpage
\seccio{sec7}{Equacions exponencials}{La incògnita està en l'exponent.}

\vfill
\begin{tcolorbox}[enhanced,colback=fmyellow!18,colframe=fmgold,arc=4mm,boxrule=1.6pt,before upper=\centering,top=2mm,bottom=2mm]
{\Large$a^M=a^N\iff M=N$}\par\smallskip
{\small $a>0$, $a\neq1$}
\end{tcolorbox}

\vfill
\begin{tcbraster}[raster columns=2,raster equal height=rows,raster column skip=10pt,raster row skip=8pt]
\begin{tcolorbox}[fmbase,colback=fmK!8,colframe=fmK,coltitle=white,title={\pas{1}\ Mateixa base}]
Igualem els exponents:\par\smallskip
$2^{x-1}=16=2^4\ \to\ x-1=4\ \to\ \boxed{x=5}$
\end{tcolorbox}
\begin{tcolorbox}[fmbase,colback=fmA!8,colframe=fmA,coltitle=white,title={\pas{2}\ Bases diferents}]
Prenem \textbf{logaritmes}:\par\smallskip
$5^x=12\ \to\ x=\dfrac{\log12}{\log5}\approx\boxed{1{,}544}$
\end{tcolorbox}
\begin{tcolorbox}[fmbase,colback=fmB!8,colframe=fmB,coltitle=white,title={\pas{3}\ Sumes de potències}]
Factor comú $a^x$:\par\smallskip
$3^{x+1}+3^x+3^{x-1}=39$\\ $3^x\left(3+1+\tfrac13\right)=39\ \to\ 3^x=9\ \to\ \boxed{x=2}$
\end{tcolorbox}
\begin{tcolorbox}[fmbase,colback=fmV!8,colframe=fmV,coltitle=white,title={\pas{4}\ Apareixen $a^{2x}$ i $a^x$}]
Canvi $t=a^x$:\par\smallskip
$4^x-6\cdot2^x+8=0\ \to\ t^2-6t+8=0$\\ $t=4\to x=2$; \ $t=2\to x=1$
\end{tcolorbox}
\end{tcbraster}

\vfill
\begin{compte}[$a^x$ sempre és positiu]
En el canvi $t=a^x$, les solucions $t\le0$ \textbf{no valen}: \ $9^x+5\cdot3^x-24=0$ dona $t=3$ i $t=-8$: només $x=1$.
\end{compte}

\vfill
\begin{exemple}[Exemples]
$2^{x^2-3}=64\ \to\ x^2-3=6\ \to\ x=\pm3$ \qquad $\left(\frac12\right)^x=8\ \to\ 2^{-x}=2^3\ \to\ x=-3$\\[4pt]
$2^{x+1}+2^x+2^{x-1}=28\ \to\ 2^x\cdot\frac72=28\ \to\ x=3$ \qquad $3^x=20\ \to\ x=\dfrac{\log20}{\log3}\approx2{,}727$
\end{exemple}

\vfill
\begin{tcbraster}[raster columns=2,raster equal height,raster column skip=10pt]
\begin{exemple}[Exemple: \ $\dfrac{2^{x^2}}{2^x}=4$]
$2^{x^2-x}=2^2\ \to\ x^2-x-2=0$\\[3pt]
\boxed{x=2,\ x=-1}
\end{exemple}
\begin{exemple}[Exemple: \ $3^{2x}-4\cdot3^{x+1}+27=0$]
$3^{x+1}=3\cdot3^x$; \ $t=3^x$: \ $t^2-12t+27=0$\\[3pt]
$t=9\to x=2$; \ $t=3\to x=1$
\end{exemple}
\end{tcbraster}

\vfill
\begin{exemple}[Problemes]
500 bacteris que es dupliquen cada hora. Quan seran 64\,000? \ $500\cdot2^t=64\,000\ \to\ 2^t=128\ \to\ \boxed{t=7\text{ h}}$\\[5pt]
1000 € al 5\,\% anual. Quan seran 1500 €? \ $1{,}05^t=1{,}5\ \to\ t=\dfrac{\log1{,}5}{\log1{,}05}\approx8{,}31$ \ $\to$ \ \boxed{9\text{ anys}}
\end{exemple}

\vfill
% ============================================================ 3.8
\clearpage
\seccio{sec8}{Equacions logarítmiques}{Un logaritme a cada membre, o la definició.}

\vfill
\begin{minipage}[c]{0.48\linewidth}
\begin{idea}[Recorda]
$\log_aN=x\iff a^x=N$\\[3pt]
$\log(M\cdot N)=\log M+\log N$\\[2pt]
$\log\frac MN=\log M-\log N$\\[2pt]
$\log M^k=k\log M$ \qquad $1=\log10$
\end{idea}
\end{minipage}\hfill
\begin{minipage}[c]{0.49\linewidth}
\passos{\pas{1} Agrupem: \ $\log A=\log B\to A=B$.\\[3pt]
\pas{2} O bé $\log_aA=k\to A=a^k$.\\[3pt]
\pas{3} \textbf{Comprovem}: tots els arguments han de ser \textbf{positius}.}
\end{minipage}

\vfill
\begin{exemple}[Exemple: \ per la definició]
$\log(x+1)=2\ \to\ x+1=100\ \to\ \boxed{x=99}$ \qquad $\log_2(x-3)=3\ \to\ x-3=8\ \to\ \boxed{x=11}$
\end{exemple}

\vfill
\begin{exemple}[Exemple: \ $\log x+\log(x-3)=1$]
$\log\big(x(x-3)\big)=\log10\ \to\ x^2-3x-10=0\ \to\ x=5,\ x=-2$\\[3pt]
$\log(-2)$ no existeix: \ \mal{x=-2} \qquad Solució: \boxed{x=5}
\end{exemple}

\vfill
\begin{exemple}[Exemple: \ $2\log x-\log(x+6)=0$]
$\log x^2=\log(x+6)\ \to\ x^2-x-6=0\ \to\ x=3,\ x=-2$ \qquad només val \boxed{x=3}
\end{exemple}

\vfill
\begin{exemple}[Exemples]
$\log_3(2x+1)=2\to x=4$ \qquad $\ln x+\ln2=\ln10\to x=5$ \qquad $\log(15-2x)=2\log x\to x=3$ \ \nota{$-5$ no val}
\end{exemple}

\vfill
\begin{exemple}[Exemple: \ $\log(x+6)-\log(x-3)=1$]
$\log\dfrac{x+6}{x-3}=\log10\ \to\ x+6=10(x-3)\ \to\ 9x=36\ \to\ \boxed{x=4}$ \ \be{}
\end{exemple}

\vfill
\begin{tcbraster}[raster columns=2,raster equal height,raster column skip=10pt]
\begin{exemple}[Exemple: \ $\log_2x+\log_2(x+2)=3$]
$\log_2\big(x(x+2)\big)=3\ \to\ x^2+2x=2^3$\\[3pt]
$x^2+2x-8=0\ \to\ x=2,\ x=-4$\\[3pt]
$x=-4$ no val: \ \boxed{x=2}
\end{exemple}
\begin{exemple}[Exemple: \ $(\log x)^2-3\log x+2=0$]
Canvi $t=\log x$: \ $t^2-3t+2=0$\\[3pt]
$t=1\to x=10$; \ $t=2\to x=100$\\[3pt]
\boxed{x=10,\ x=100}
\end{exemple}
\end{tcbraster}

\vfill
\begin{minipage}[t]{0.5\linewidth}
\begin{exemple}[Exemple: \ $2\log x=\log(5x-6)$]
$x^2=5x-6\ \to\ x^2-5x+6=0$\\[3pt]
$x=2$ i $x=3$: totes dues valen \be{}
\end{exemple}
\end{minipage}\hfill
\begin{minipage}[t]{0.47\linewidth}
\begin{compte}[El logaritme d'una suma]
$\log(x+3)\neq\log x+\log3$. \ La propietat és per al \textbf{producte}: \ $\log(3x)=\log3+\log x$.
\end{compte}
\end{minipage}

\vfill
% ============================================================ 3.9
\clearpage
\seccio{sec9}{Inequacions de primer grau}{Com una equació, però amb una regla especial.}

\vfill
\begin{idea}[Inequació]
Una desigualtat ($<$, $\le$, $>$, $\ge$) amb una incògnita. Les solucions solen ser \textbf{infinites}: les escrivim amb \textbf{intervals}.
\end{idea}

\vfill
\begin{tcbraster}[raster columns=2,raster equal height,raster column skip=10pt]
\begin{tcolorbox}[fmbase,colback=fmK!8,colframe=fmK,coltitle=white,title={\faCheck\ \ Es manté}]
Sumar o restar, i multiplicar o dividir per un nombre \textbf{positiu}.
\end{tcolorbox}
\begin{tcolorbox}[fmbase,colback=fmR!6,colframe=fmR,coltitle=white,title={\faIcon{exchange-alt}\ \ Canvia de sentit}]
Multiplicar o dividir per un \textbf{negatiu}:\par\smallskip
{\centering$-2x<6\ \to\ x\cR{>}-3$\par}
\end{tcolorbox}
\end{tcbraster}

\vfill
\begin{tcbraster}[raster columns=2,raster equal height,raster column skip=10pt]
\begin{exemple}[Exemple: \ $3x-2\le x+4$]
$2x\le6\ \to\ x\le3$ \qquad $\boxed{\lo-\infty,\,3]}$
\end{exemple}
\begin{exemple}[Exemple: \ $-2x+1<7$]
$-2x<6$. \ Entre $-2$, \textbf{girem}: \ $x>-3$ \qquad $\boxed{\lo-3,+\infty\ro}$
\end{exemple}
\end{tcbraster}

\vfill
\begin{exemple}[Exemple: \ $\dfrac{x-1}{2}-\dfrac{x+2}{3}\ge1$]
Per 6 (positiu): \ $3(x-1)-2(x+2)\ge6\ \to\ 3x-3-2x-4\ge6\ \to\ x\ge13$ \qquad $\boxed{[13,+\infty\ro}$
\end{exemple}

\vfill
\begin{exemple}[Exemples]
$5x+3>2x-6\to x>-3$ \qquad $4-3x\ge10\to x\le-2$ \qquad $2(x-1)<5x+4\to x>-2$ \qquad $\frac x2-1\le\frac x3\to x\le6$\\[4pt]
Amb parèntesis: \ $3(x-2)-2(1-x)>4x+1\ \to\ 5x-8>4x+1\ \to\ x>9$: \ $\lo9,+\infty\ro$
\end{exemple}

\vfill
\begin{exemple}[Exemple: \ sistema \ $2x-1>3$ \ i \ $x+4\le10$]
\begin{minipage}[c]{0.5\linewidth}
Resolem cada una: \ $x>2$ \ i \ $x\le6$\\[3pt]
Les solucions són les \textbf{comunes} (intersecció):\\[3pt]
$\boxed{\lo2,\,6]}$
\end{minipage}\hfill
\begin{minipage}[c]{0.47\linewidth}\centering
\begin{tikzpicture}[x=0.5cm]
  \draw[fmmuted,line width=0.8pt] (-1,0) -- (9,0); \foreach \k in {0,...,8} \draw[fmmuted] (\k,-0.07) -- (\k,0.07);
  \draw[fmA,line width=2.4pt,-{Latex[length=2.4mm]}] (2,0.35) -- (9,0.35); \filldraw[fill=white,draw=fmA,line width=1.2pt] (2,0.35) circle (3pt);
  \draw[fmB,line width=2.4pt,{Latex[length=2.4mm]}-] (-1,-0.35) -- (6,-0.35); \fill[fmB] (6,-0.35) circle (3.2pt);
  \fill[fmK!40,opacity=0.6] (2,-0.15) rectangle (6,0.15);
  \node[below=6pt,font=\scriptsize] at (2,-0.35) {2}; \node[below=6pt,font=\scriptsize] at (6,-0.35) {6};
\end{tikzpicture}
\end{minipage}
\end{exemple}

\vfill
\begin{minipage}[t]{0.48\linewidth}
\begin{exemple}[Exemple: \ $-3<2x+1\le7$]
Restem 1 als tres membres: \ $-4<2x\le6$\\[3pt]
Dividim entre 2: \ $-2<x\le3$ \ $\to$ \ $\boxed{\lo-2,\,3]}$
\end{exemple}
\end{minipage}\hfill
\begin{minipage}[t]{0.49\linewidth}
\begin{exemple}[Problema: \ tarifa A: 20 € + 0,05 €/min; \ tarifa B: 0,10 €/min. Quan és millor A?]
$20+0{,}05x<0{,}10x\ \to\ 20<0{,}05x$\\[3pt]
$x>400$: \ \boxed{\text{més de 400 min}}
\end{exemple}
\end{minipage}

\vfill
% ============================================================ 3.10
\clearpage
\seccio{sec10}{Inequacions de segon grau o superior}{Arrels, intervals i taula de signes.}

\vfill
\begin{minipage}[c]{0.52\linewidth}
\passos{\pas{1} Tot a un membre: \ $P(x)>0$ \ o \ $P(x)<0$ \ (o amb $\ge$, $\le$).\\[3pt]
\pas{2} Resolem l'equació $P(x)=0$.\\[3pt]
\pas{3} Situem les solucions en la recta real i vegem el signe en cada tros substituint en $P(x)$.\\[3pt]
\pas{4} Triem els intervals. Les arrels, només amb $\le$ o $\ge$.}
\end{minipage}\hfill
\begin{minipage}[c]{0.45\linewidth}\centering
\begin{tikzpicture}[x=0.7cm,y=0.55cm]
  \draw[-{Latex[length=2mm]},fmmuted,line width=0.8pt] (-0.5,0) -- (5.5,0) node[right,font=\scriptsize] {$x$};
  \draw[fmblue,line width=1.6pt,domain=0.9:4.1,samples=40] plot (\x,{(\x-2)*(\x-3)*1.5});
  \draw[fmK,line width=4pt,opacity=0.6] (-0.5,0) -- (2,0) (3,0) -- (5.5,0);
  \draw[fmR,line width=4pt,opacity=0.6] (2,0) -- (3,0);
  \node[below,font=\scriptsize] at (2,-0.1) {2}; \node[below,font=\scriptsize] at (3,-0.1) {3};
  \node[font=\small\heavy,text=fmK!80!black] at (0.8,1.2) {$+$}; \node[font=\small\heavy,text=fmR] at (2.5,-1) {$-$}; \node[font=\small\heavy,text=fmK!80!black] at (4.2,1.2) {$+$};
\end{tikzpicture}\\
{\comic\small\color{fmmuted}$a>0$: negativa \textbf{entre} les arrels}
\end{minipage}


\vfill
\begin{exemple}[Exemple: \ $x^3-x^2-6x\ge0$]
\begin{minipage}[c]{0.62\linewidth}\centering\small
\renewcommand{\arraystretch}{1.25}
\taularod[fmex]{>{\heavy}c*{7}{>{\hspace{2pt}}c<{\hspace{2pt}}}}{
\rowcolor{fmex} & \color{white}$\lo-\infty,-2\ro$ & \color{white}$-2$ & \color{white}$\lo-2,0\ro$ & \color{white}$0$ & \color{white}$\lo0,3\ro$ & \color{white}$3$ & \color{white}$\lo3,+\infty\ro$\\
$x+2$ & $-$ & 0 & $+$ & $+$ & $+$ & $+$ & $+$\\
\rowcolor{fmexbg}$x$ & $-$ & $-$ & $-$ & 0 & $+$ & $+$ & $+$\\
$x-3$ & $-$ & $-$ & $-$ & $-$ & $-$ & 0 & $+$\\
\rowcolor{fmexbg}$P(x)$ & \cR{$-$} & 0 & \cK{$+$} & 0 & \cR{$-$} & 0 & \cK{$+$}
}
\end{minipage}\hfill
\begin{minipage}[c]{0.35\linewidth}
$x(x+2)(x-3)\ge0$\\[4pt]
$\boxed{[-2,\,0]\cup[3,+\infty\ro}$
\end{minipage}
\end{exemple}

\vfill
\begin{exemple}[Exemples]
$-x^2+x+2\ge0\ \to\ x^2-x-2\le0$ \ \nota{per $-1$, girem}: \ $[-1,\,2]$ \qquad $x^2\ge9$: \ $\lo-\infty,-3]\cup[3,+\infty\ro$\\[4pt]
$(x+1)^2>0$: \ $\mathbb{R}-\{-1\}$ \qquad $x^2+1<0$: \ \textbf{cap solució} \qquad $x^3-4x>0$: \ $\lo-2,0\ro\cup\lo2,+\infty\ro$
\end{exemple}

\vfill
\begin{exemple}[Exemple: \ $x^2-x-6\le0$, \ amb la recta]
\begin{minipage}[c]{0.45\linewidth}
\pas{1} $x^2-x-6=0\ \to\ x=-2,\ x=3$\\[3pt]
\pas{2} Un valor de cada tros:\\ $P(-3)=6$, \ $P(0)=-6$, \ $P(4)=6$\\[3pt]
\pas{3} Volem $\le0$: \ $\boxed{[-2,\,3]}$ \ \nota{amb les arrels}
\end{minipage}\hfill
\begin{minipage}[c]{0.52\linewidth}\centering
\begin{tikzpicture}[x=0.75cm]
  \draw[-{Latex[length=2mm]},fmmuted,line width=0.8pt] (-5,0) -- (5.5,0);
  \draw[fmR,line width=4pt,opacity=0.6] (-2,0) -- (3,0);
  \fill[fmnavy] (-2,0) circle (3pt) node[below=3pt,font=\scriptsize] {$-2$}; \fill[fmnavy] (3,0) circle (3pt) node[below=3pt,font=\scriptsize] {$3$};
  \node[font=\heavy,text=fmK!80!black] at (-3.6,0.45) {$+$}; \node[font=\heavy,text=fmR] at (0.5,0.45) {$-$}; \node[font=\heavy,text=fmK!80!black] at (4.3,0.45) {$+$};
  \node[font=\scriptsize,text=fmmuted] at (-3.6,-0.45) {$x=-3$}; \node[font=\scriptsize,text=fmmuted] at (0.5,-0.45) {$x=0$}; \node[font=\scriptsize,text=fmmuted] at (4.3,-0.45) {$x=4$};
\end{tikzpicture}
\end{minipage}
\end{exemple}

\vfill
\begin{tcbraster}[raster columns=2,raster equal height,raster column skip=10pt]
\begin{exemple}[Exemple: \ $(x+1)(x-3)<2x-3$]
$x^2-2x-3<2x-3\ \to\ x^2-4x<0$\\[3pt]
$x(x-4)<0$: \ arrels 0 i 4; \ $P(1)=-3<0$\\[3pt]
$\boxed{\lo0,\,4\ro}$
\end{exemple}
\begin{exemple}[Problema: \ rectangle de 20 m de perímetre. Quina base fa que l'àrea passe de 21 m$^2$?]
$x(10-x)>21\ \to\ x^2-10x+21<0$\\[3pt]
Arrels 3 i 7: \ \boxed{\text{entre 3 i 7 m}}
\end{exemple}
\end{tcbraster}

\vfill
% ============================================================ resum
\resumtitol
\begin{tcbraster}[raster columns=2,raster equal height=rows,raster column skip=14pt,raster row skip=30pt,raster before skip=16pt]
\rtile{3.1}{Segon grau}{$x=\frac{-b\pm\sqrt{\Delta}}{2a}$. \ Incompletes: factor comú o aïllar $x^2$.}
\rtile{3.2}{Biquadrades}{$t=x^2$; \ $t<0$ no dona solucions.}
\rtile{3.3}{Grau $>2$}{Ruffini amb els divisors del terme independent.}
\rtile{3.4}{Factoritzades}{Cada factor igual a 0 (si l'altre membre és 0).}
\rtile{3.5}{Fraccions}{Pel m.c.m. i comprovem els denominadors.}
\rtile{3.6}{Radicals}{Aïllem, al quadrat i comprovem.}
\rtile{3.7}{Exponencials}{Mateixa base, logaritmes, factor comú o $t=a^x$.}
\rtile{3.8}{Logarítmiques}{Un log a cada membre o la definició. Arguments positius.}
\rtile{3.9}{Inequacions 1r grau}{Per un negatiu, la desigualtat gira.}
\rtile{3.10}{Grau 2 o més}{Arrels i taula de signes.}
\end{tcbraster}
\end{document}
